% Hartshorne, Algebraic Geometry (1977), Chapter II.
% Source-faithful mathematical blueprint of Hartshorne, Algebraic Geometry (1977), Chapter II.
% Authorial exposition, examples, and exercises are omitted; definitions, statements, and proof
% arguments are retained.

\section{Sheaves}

\begin{definition}[Presheaf and sheaf]
\label{def:II-1-presheaf}
\source{hartshorne-algebraic-geometry:II.1}
\lean{Hartshorne.Presheaf, Hartshorne.Sheaf, Hartshorne.IsSheaf}
\leanok
A presheaf $\mathcal F$ on $X$ assigns an object $\mathcal F(U)$ to every open $U$ and restriction maps
$\mathcal F(U)\to\mathcal F(V)$ for $V\subset U$, functorially, with $\mathcal F(\varnothing)=0$.
It is a sheaf when sections that agree locally glue uniquely. The definitions apply to sets,
abelian groups, rings, and modules.
More explicitly, if $\{U_i\}$ covers $U$, then restriction to the $U_i$ is injective, and every
family $s_i\in\mathcal F(U_i)$ satisfying
$s_i|_{U_i\cap U_j}=s_j|_{U_i\cap U_j}$ comes from a unique $s\in\mathcal F(U)$.
\end{definition}

\begin{definition}[Stalks and morphisms]
\label{def:II-1-stalk}
\source{hartshorne-algebraic-geometry:II.1}
\lean{Hartshorne.Stalk, Hartshorne.germ, Hartshorne.stalkMap, Hartshorne.germ_restrict}
\leanok
The stalk is $\mathcal F_x=\varinjlim_{x\in U}\mathcal F(U)$. A morphism is a compatible family
of maps on opens and induces maps on stalks.  A germ is represented by a pair $(U,s)$ with $x\in U$;
$(U,s)$ and $(V,t)$ define the same germ exactly when $s$ and $t$ agree on some neighborhood of $x$
contained in $U\cap V$.
\uses{def:II-1-presheaf}
\end{definition}

\begin{proposition}[Sheaf isomorphisms are stalkwise]
\label{prop:II-1-1}
\source{hartshorne-algebraic-geometry:II.1.1}
\dcref{II.1.1}
\lean{Hartshorne.sheaf_iso_iff_stalkMap_iso}
\leanok
A morphism of sheaves is an isomorphism if and only if it induces an isomorphism on every stalk.
\uses{def:II-1-presheaf,def:II-1-stalk}
\end{proposition}
\begin{proof}
The forward implication is immediate. If all stalk maps are injective, a section whose germs vanish
has a cover on which it vanishes, hence is zero. If all are surjective, lift the germ of a target
section near each point; after shrinking, the lifts agree on overlaps by injectivity, and the sheaf
axiom glues them. The glued lift maps to the target section.
\end{proof}


\begin{proposition}[Sheafification]
\label{prop:II-1-2}
\source{hartshorne-algebraic-geometry:II.1.2}
\dcref{II.1.2}
\lean{Hartshorne.sheafification, Hartshorne.sheafification_condition,
Hartshorne.sheafification_stalkIso, Hartshorne.categoricalSheafification,
Hartshorne.categoricalSheafificationUnit, Hartshorne.categoricalSheafificationLift,
Hartshorne.categoricalSheafificationUnit_lift,
Hartshorne.categoricalSheafificationLift_unique,
Hartshorne.categoricalSheafificationUnit_stalk_isIso}
\leanok
Every presheaf $\mathcal F$ admits a sheafification $\mathcal F^a$ and a universal map
$\mathcal F\to\mathcal F^a$; it is unique up to unique isomorphism and has the same stalks as
$\mathcal F$.
\uses{def:II-1-presheaf,def:II-1-stalk}
\end{proposition}
\begin{proof}
Let $\mathcal F^a(U)$ be the functions assigning to each $x\in U$ a germ in $\mathcal F_x$ which,
locally at every $x$, is represented by a section of $\mathcal F$. These functions form a sheaf under
restriction; the germ map is universal because any map to a sheaf is determined locally.
\end{proof}

\begin{definition}[Direct and inverse image]
\label{def:II-1-direct-inverse}
\source{hartshorne-algebraic-geometry:II.1}
\lean{Hartshorne.directImage, Hartshorne.presheafDirectImage,
Hartshorne.presheafInverseImage, Hartshorne.sheafInverseImage,
Hartshorne.sheafInverseImageDirectImageAdjunction,
Hartshorne.directImage_presheaf_isSheaf}
\leanok
For a continuous $f:X\to Y$, define $(f_*\mathcal F)(V)=\mathcal F(f^{-1}V)$. The inverse image
$f^{-1}\mathcal G$ is the sheafification of $U\mapsto\varinjlim_{f(U)\subset V}\mathcal G(V)$.
Restriction to a subspace is inverse image under its inclusion.
\uses{def:II-1-presheaf,prop:II-1-2}
\end{definition}

For a morphism $\varphi:\mathcal F\to\mathcal G$ of presheaves, define the presheaf kernel,
cokernel, and image by
\[
U\longmapsto\ker\varphi(U),\qquad U\longmapsto\operatorname{coker}\varphi(U),
\qquad U\longmapsto\operatorname{im}\varphi(U).
\]
The kernel presheaf of a morphism of sheaves is a sheaf; in general the presheaf cokernel and image
must be sheafified.  Thus kernels are computed sectionwise, while cokernels and images in the sheaf
category are the associated sheaves of the sectionwise constructions.

\begin{definition}[Subsheaves, exactness, and quotients]
\label{def:II-1-subsheaf-exact}
\source{hartshorne-algebraic-geometry:II.1}
A subsheaf $\mathcal F'\subseteq\mathcal F$ is a sheaf whose sections are subgroups of the
sections of $\mathcal F$ and whose restrictions are induced.  The kernel of a sheaf morphism is its
sectionwise kernel, the image is the sheaf associated to its sectionwise image, and a morphism is
injective or surjective according as its kernel is zero or its image is the target.  A sequence is
exact when the image at every term equals the kernel at that term.  The quotient $\mathcal F/\mathcal F'$
is the sheaf associated to $U\mapsto\mathcal F(U)/\mathcal F'(U)$, with stalks
$(\mathcal F/\mathcal F')_x=\mathcal F_x/\mathcal F'_x$.
\uses{def:II-1-presheaf,def:II-1-stalk,prop:II-1-2}
\end{definition}

\section{Schemes}

\begin{lemma}[Closed sets in a spectrum]
\label{lem:II-2-1}
\source{hartshorne-algebraic-geometry:II.2.1}
\dcref{II.2.1}
\lean{Hartshorne.spectrumZeroLocus_mul, Hartshorne.spectrumZeroLocus_iSup,
Hartshorne.spectrumZeroLocus_eq_iff}
\leanok
For ideals $\mathfrak a,\mathfrak b\subset A$,
$V(\mathfrak a\mathfrak b)=V(\mathfrak a)\cup V(\mathfrak b)$,
$V(\sum_i\mathfrak a_i)=\bigcap_iV(\mathfrak a_i)$, and
$V(\mathfrak a)=V(\mathfrak b)$ iff $\sqrt{\mathfrak a}=\sqrt{\mathfrak b}$.
\end{lemma}
\begin{proof}
Use primality for the product formula, the definition of generated ideals for intersections, and
the fact that the radical is the intersection of all primes containing an ideal.
\end{proof}

\begin{definition}[Affine spectrum]
\label{def:II-2-spectrum}
\source{hartshorne-algebraic-geometry:II.2}
\lean{Hartshorne.AffineSpectrum, Hartshorne.spectrumZeroLocus,
Hartshorne.spectrumBasicOpen, Hartshorne.spectrumBasicOpen_isTopologicalBasis,
Hartshorne.affineSpec, Hartshorne.affineStructureSheaf}
 $\operatorname{Spec}A$ is the set of prime ideals with closed sets $V(\mathfrak a)$. Its structure
sheaf consists of locally represented fractions in the local rings $A_{\mathfrak p}$; $D(f)$ is a
basis of opens.
\uses{lem:II-2-1,def:II-1-presheaf}
\end{definition}

\begin{proposition}[Sections of the affine structure sheaf]
\label{prop:II-2-2}
\source{hartshorne-algebraic-geometry:II.2.2}
\dcref{II.2.2}
\lean{Hartshorne.affineStructureSheaf_stalk_isLocalization,
Hartshorne.affineStructureSheaf_stalk_iso, Hartshorne.affineSpec_stalk_iso,
Hartshorne.affineStructureSheaf_basicOpen_isLocalization,
Hartshorne.affineStructureSheaf_basicOpen_iso,
Hartshorne.affineStructureSheaf_globalSections_bijective,
Hartshorne.affineStructureSheaf_globalSectionsIso,
Hartshorne.affineSpec_globalSectionsIso}
\leanok
For $X=\operatorname{Spec}A$,
$\mathcal O_{X,\mathfrak p}=A_{\mathfrak p},\qquad
\Gamma(D(f),\mathcal O_X)=A_f,\qquad\Gamma(X,\mathcal O_X)=A$.
\uses{def:II-2-spectrum,def:II-1-stalk}
\end{proposition}
\begin{proof}
Evaluation gives the stalk map to $A_{\mathfrak p}$; a fraction represents a section on $D(f)$, proving
surjectivity. Equality of germs is equality after multiplying by an element outside $\mathfrak p$,
which proves injectivity. For $D(f)$, injectivity follows by taking annihilators and surjectivity by a
finite $D(h_i)$-cover, clearing denominators on overlaps, and using a power of $f$ in $(h_i)$.
\end{proof}

\begin{definition}[Ringed and locally ringed spaces]
\label{def:II-2-ringed}
\source{hartshorne-algebraic-geometry:II.2}
A ringed space is $(X,\mathcal O_X)$; a morphism $(f,f^\sharp)$ has a continuous $f$ and
$f^\sharp:\mathcal O_Y\to f_*\mathcal O_X$. It is locally ringed when all stalks are local and
morphisms induce local homomorphisms.
\uses{def:II-1-presheaf,def:II-1-stalk,def:II-1-direct-inverse}
\end{definition}

\begin{proposition}[Affine schemes and ring maps]
\label{prop:II-2-3}
\source{hartshorne-algebraic-geometry:II.2.3}
\dcref{II.2.3}
\lean{Hartshorne.affineSpecMap, Hartshorne.affineSpecMap_apply,
Hartshorne.affineSpecMap_continuous, Hartshorne.affineSpec_homEquiv,
Hartshorne.affineSpec_homEquiv_map, Hartshorne.affineSpec_map_homEquiv,
Hartshorne.affineSpec_stalk_isLocalRing,
Hartshorne.affineSpecMap_stalkMap_isLocalHom}
\leanok
$\operatorname{Spec}A$ is locally ringed. A ring map $A\to B$ induces
$\operatorname{Spec}B\to\operatorname{Spec}A$, and every morphism
$\operatorname{Spec}B\to\operatorname{Spec}A$ arises uniquely in this way.
\uses{prop:II-2-2}
\end{proposition}
\begin{proof}
The stalk assertion is Proposition \ref{prop:II-2-2}. A ring map sends a prime to its inverse image,
is continuous on $V(\mathfrak a)$, and localizes to maps $A_{\mathfrak p}\to B_{\mathfrak q}$. Conversely,
take global sections of a locally ringed-space morphism; the induced stalk map shows its underlying
point map is inverse-image of primes, and localization compatibility determines the sheaf map.
\end{proof}

\begin{definition}[Scheme and subschemes]
\label{def:II-2-scheme}
\source{hartshorne-algebraic-geometry:II.2}
An affine scheme is a locally ringed space isomorphic to an affine spectrum; a scheme is locally
affine. Open subschemes have restricted structure sheaves. A closed immersion is a homeomorphism
onto a closed subset with surjective structure-sheaf map.
\uses{def:II-2-spectrum,def:II-2-ringed}
\end{definition}

\begin{lemma}[Gluing schemes]
\label{lem:II-2-4}
\source{hartshorne-algebraic-geometry:II.2.4}
\dcref{II.2.4}
Schemes and morphisms glue from affine open charts and compatible overlap isomorphisms satisfying
the cocycle condition.
\end{lemma}
\begin{proof}
Glue the topological spaces and glue the compatible structure sheaves using the sheaf axioms;
the affine charts descend because overlaps are open subschemes. Morphisms glue locally.
\end{proof}

\begin{definition}[Proj]
\label{def:II-2-proj}
\source{hartshorne-algebraic-geometry:II.2}
For a graded ring $S$, $\operatorname{Proj}S$ is the homogeneous-prime space excluding the irrelevant
ideal, with basic opens $D_+(f)$, and structure sheaf given by degree-zero localizations.
\uses{def:II-2-spectrum,def:II-1-presheaf}
\end{definition}

\begin{proposition}[Proj charts]
\label{prop:II-2-5}
\source{hartshorne-algebraic-geometry:II.2.5}
\dcref{II.2.5}
The $D_+(f)$ form an affine open cover of $\operatorname{Proj}S$ and
$\Gamma(D_+(f),\mathcal O)=S_{(f)}$.
\uses{prop:II-2-2}
\end{proposition}
\begin{proof}
The homogeneous-prime topology is covered by $D_+(f)$ for positive-degree homogeneous $f$. Degree-zero
fractions glue exactly as in Proposition \ref{prop:II-2-2}, giving the affine chart ring $S_{(f)}$.
\end{proof}

\begin{proposition}[Projective space represents line bundles]
\label{prop:II-2-6}
\source{hartshorne-algebraic-geometry:II.2.6}
\dcref{II.2.6}
For a field $k$, a morphism $X\to\mathbb P^n_k$ is equivalent to an invertible sheaf $L$ on $X$ and
global sections $s_0,\ldots,s_n$ generating $L$, up to simultaneous multiplication by an isomorphism.
\uses{def:II-2-scheme,def:II-1-subsheaf-exact}
\end{proposition}
\begin{proof}
Pull back $\mathcal O(1)$ and its coordinate sections. Conversely, the ratios $s_i/s_j$ define compatible
maps to the standard affine charts; gluing gives the morphism and the two constructions are inverse.
\end{proof}

\section{First Properties of Schemes}

\begin{definition}[Connected, reduced, integral, normal]
\label{def:II-3-properties}
\source{hartshorne-algebraic-geometry:II.3}
A scheme is connected or irreducible according to its topological space, reduced when all affine
rings are reduced, integral when reduced and irreducible, and locally Noetherian when it has a cover
by spectra of Noetherian rings. An integral scheme is normal if its local rings are integrally closed.
\end{definition}

\begin{proposition}[Integral schemes]
\label{prop:II-3-1}
\source{hartshorne-algebraic-geometry:II.3.1}
\dcref{II.3.1}
A scheme is integral iff it is reduced and irreducible.
\end{proposition}
\begin{proof}
On an affine open, irreducibility says the nilradical is prime; reducedness makes it zero, so the
coordinate ring is a domain. The converse follows because spectra of domains are irreducible and
localization preserves domains.
\end{proof}

\begin{proposition}[Noetherian schemes]
\label{prop:II-3-2}
\source{hartshorne-algebraic-geometry:II.3.2}
\dcref{II.3.2}
A scheme is locally Noetherian iff every affine open has Noetherian coordinate ring; an affine scheme
is Noetherian iff its coordinate ring is Noetherian.
\end{proposition}
\begin{proof}
Localizations of Noetherian rings are Noetherian, so noetherian affine charts form a basis. If
$D(f_i)$ cover $\operatorname{Spec}A$ and each $A_{f_i}$ is Noetherian, quasi-compactness gives finitely
many $f_i$ generating the unit ideal. For an ideal chain, localization makes each chain stationary;
the equality $\mathfrak a=\bigcap_i(\mathfrak a A_{f_i}\cap A)$ proves the original chain stationary.
\end{proof}

\begin{definition}[Morphisms of finite type and finite]
\label{def:II-3-finite-type}
\source{hartshorne-algebraic-geometry:II.3}
A morphism is locally of finite type when affine chart rings are finitely generated algebras, of
finite type when finitely many charts suffice over each affine target, and finite when inverse images
of affine opens are affine with finite module coordinate rings.
\end{definition}

\begin{definition}[Dimension and codimension]
\label{def:II-3-dimension}
\source{hartshorne-algebraic-geometry:II.3}
Dimension is the supremum of chains of irreducible closed subsets; codimension is the supremum of
chains beginning at a given irreducible closed subset. For $\operatorname{Spec}A$ this is Krull dimension.
\end{definition}

\begin{theorem}[Fibred products]
\label{thm:II-3-3}
\source{hartshorne-algebraic-geometry:II.3.3}
\dcref{II.3.3}
For schemes $X,Y$ over $S$, $X\times_S Y$ exists and is unique up to unique isomorphism; for affine
schemes it is $\operatorname{Spec}(A\otimes_RB)$.
\uses{prop:II-2-3}
\end{theorem}
\begin{proof}
The affine case is the tensor-product universal property and Proposition \ref{prop:II-2-3}.
Products over affine opens are glued along products over overlaps. The resulting projections and
universal morphism glue from the corresponding affine constructions; uniqueness is local.
\end{proof}

\begin{definition}[Fibres and base extension]
\label{def:II-3-fibre}
\source{hartshorne-algebraic-geometry:II.3}
For $f:X\to Y$ and $y\in Y$, $X_y=X\times_Y\operatorname{Spec}\kappa(y)$. Base extension is
$X_{S'}=X\times_SS'$. Fibres and base extensions are associative up to canonical isomorphism.
\end{definition}

\section{Separated and Proper Morphisms}

\begin{definition}[Separated morphism]
\label{def:II-4-separated}
\source{hartshorne-algebraic-geometry:II.4}
A morphism $f:X\to Y$ is separated when its diagonal $\Delta_f:X\to X\times_YX$ is a closed immersion.
\end{definition}

\begin{proposition}[Affine morphisms are separated]
\label{prop:II-4-1}
\source{hartshorne-algebraic-geometry:II.4.1}
\dcref{II.4.1}
Every morphism of affine schemes is separated.
\uses{prop:II-2-3}
\end{proposition}
\begin{proof}
For $A\to B$, the diagonal corresponds to the multiplication map $B\otimes_AB\to B$, which is surjective;
Proposition \ref{prop:II-2-3} identifies it with a closed immersion.
\end{proof}

\begin{corollary}[Separatedness is local on the target]
\label{cor:II-4-2}
\source{hartshorne-algebraic-geometry:II.4.2}
\dcref{II.4.2}
A morphism is separated iff its restriction over every member of an open affine cover of the target is
separated.
\end{corollary}
\begin{proof}
The diagonal is compatible with base change by target opens; closed immersions are local on the target.
\end{proof}

\begin{theorem}[Valuative criterion for separatedness]
\label{thm:II-4-3}
\source{hartshorne-algebraic-geometry:II.4.3}
\dcref{II.4.3}
For a finite-type morphism $f:X\to Y$, $f$ is separated iff for every valuation ring $R$ with fraction field $K$,
a diagram $\operatorname{Spec}K\to X$ over $\operatorname{Spec}R\to Y$ has at most one extension
$\operatorname{Spec}R\to X$.
\end{theorem}
\begin{proof}
Two extensions define a map to the fibre product agreeing on the generic point, hence a map into the
diagonal. A closed diagonal contains the generic point of the valuation spectrum only through its
closed point, proving uniqueness. Conversely, failure of closedness gives a specialization in the
complement; the affine localization and domination lemma produce a valuation ring realizing two
extensions.
\end{proof}

\begin{lemma}[Valuation-ring diagrams]
\label{lem:II-4-4}
\source{hartshorne-algebraic-geometry:II.4.4}
\dcref{II.4.4}
Let $R$ be a valuation ring with fraction field $K$, $T=\operatorname{Spec}R$, and
$U=\operatorname{Spec}K$. A map $U\to X$ is a point $x_\eta$ with $k(x_\eta)=K$; a map
$T\to X$ is points $x_\eta,x_0$ with $x_0$ a specialization of $x_\eta$ and an inclusion
$k(x_\eta)=K$ such that $R$ dominates the local ring of $x_0$ on its reduced irreducible closure.
\end{lemma}
\begin{proof}
The generic point gives the first description. A map from $T$ factors through the reduced closure
of the image of its generic point, and the local homomorphism at the closed point is exactly a
domination by $R$; the converse is obtained by composing $\operatorname{Spec}R\to
\operatorname{Spec}\mathcal O_{x_0}$ with the inclusion into $X$.
\end{proof}

\begin{lemma}[Quasi-compact image criterion]
\label{lem:II-4-5}
\source{hartshorne-algebraic-geometry:II.4.5}
\dcref{II.4.5}
For a quasi-compact morphism $f:X\to Y$, the image $f(X)$ is closed if and only if it is stable
under specialization.
\end{lemma}
\begin{proof}
For the nontrivial implication, reduce to reduced affine $X,Y$ and write $X$ as a finite union of
affines. For a point in the closure of the image, choose a minimal prime below its corresponding
prime; localization and lying-over produce a point of one affine piece over that minimal point.
Stability under specialization then puts the original point in the image.
\end{proof}

\begin{corollary}[Stability of separated morphisms]
\label{cor:II-4-6}
\source{hartshorne-algebraic-geometry:II.4.6}
\dcref{II.4.6}
For noetherian schemes: open and closed immersions are separated; compositions, products over a
base, and base changes of separated morphisms are separated; if $gf$ is separated then $f$ is
separated when $g$ is separated; and separatedness is local on the target.
\end{corollary}
\begin{proof}
Apply the valuative criterion for separatedness and the fibre-product description of base change
and products; the diagonal criterion is local on the target.
\end{proof}

\begin{definition}[Proper morphism]
\label{def:II-4-proper}
\source{hartshorne-algebraic-geometry:II.4}
A morphism is proper if it is separated, of finite type, and universally closed.
\end{definition}

\begin{theorem}[Valuative criterion for properness]
\label{thm:II-4-7}
\source{hartshorne-algebraic-geometry:II.4.7}
\dcref{II.4.7}
For a finite-type morphism f, properness is equivalent to existence and uniqueness of a lift in every
valuation-ring diagram as above.
\end{theorem}
\begin{proof}
Properness gives existence because the image of the generic point in the closure of the image of
$\operatorname{Spec}R$ is closed and affine-localization yields a dominating valuation ring; uniqueness
is separatedness. Conversely, apply the lifting property to valuation rings dominating local rings
at points of a closure to show every image is closed, then use the separatedness criterion.
\end{proof}

\begin{corollary}[Stability of properness]
\label{cor:II-4-8}
\source{hartshorne-algebraic-geometry:II.4.8}
\dcref{II.4.8}
Proper morphisms are stable under composition and arbitrary base change; finite morphisms are proper.
\end{corollary}
\begin{proof}
The valuative criterion is preserved under base change and composition. A finite morphism is affine,
hence separated and finite type, and its image is closed because integral extensions satisfy lying over.
\end{proof}

\begin{theorem}[Projective morphisms are proper]
\label{thm:II-4-9}
\source{hartshorne-algebraic-geometry:II.4.9}
\dcref{II.4.9}
A projective morphism of Noetherian schemes is proper.
\end{theorem}
\begin{proof}
Closed immersions are proper. For $\mathbb P^n_Y\to Y$, homogenize the valuation-ring lifting problem
and, after scaling coordinates by a common element of K, choose a tuple in R with a unit coordinate;
this gives the lift. Apply stability under composition and closed immersion.
\end{proof}

\begin{proposition}[Varieties as schemes]
\label{prop:II-4-10}
\source{hartshorne-algebraic-geometry:II.4.10}
\dcref{II.4.10}
Over an algebraically closed field, the functor from varieties of Chapter I to schemes has image
the quasi-projective integral schemes of finite type; projective varieties correspond to projective
integral schemes. In particular the associated scheme of a variety is integral, separated, and of
finite type.
\end{proposition}
\begin{proof}
The construction on affine charts gives an integral finite-type scheme and glues to a
quasi-projective scheme. Conversely, the closed points of a projective integral scheme form an
irreducible projective variety dense in the scheme; reduced closed subschemes with the same support
are uniquely identified, and the quasi-projective case follows by an open immersion.
\end{proof}

\begin{definition}[Abstract variety]
\label{def:II-4-10-variety}
\source{hartshorne-algebraic-geometry:II.4.10}
An abstract variety is an integral separated scheme of finite type over an algebraically closed
field. A proper abstract variety is complete.
\end{definition}

\begin{theorem}[Integral closure by valuation rings]
\label{thm:II-4-11A}
\source{hartshorne-algebraic-geometry:II.4.11A}
\dcref{II.4.11A}
If $A$ is a subring of a field $K$, the integral closure of $A$ in $K$ is the intersection of all
valuation rings of $K$ containing $A$.
\end{theorem}
\begin{proof}
Hartshorne cites Bourbaki for this algebraic result and gives no proof.
\end{proof}

\section{Sheaves of Modules}

\begin{definition}[Modules and associated sheaves]
\label{def:II-5-module-sheaf}
\source{hartshorne-algebraic-geometry:II.5}
For an $A$-module $M$, define $\widetilde M$ on $\operatorname{Spec}A$ by
$\widetilde M(D(f))=M_f$. An $\mathcal O_X$-module is quasi-coherent if it is locally of this form, and
coherent on a Noetherian scheme if the modules are finitely generated.
\end{definition}

\begin{proposition}[Basic affine module identities]
\label{prop:II-5-1}
\source{hartshorne-algebraic-geometry:II.5.1}
\dcref{II.5.1}
For $\operatorname{Spec}A$, $(\widetilde M)_\mathfrak p=M_\mathfrak p$ and
$\Gamma(D(f),\widetilde M)=M_f$. The map $M\to\Gamma(\operatorname{Spec}A,\widetilde M)$ is an isomorphism.
\uses{prop:II-2-2}
\end{proposition}
\begin{proof}
The proof is the localization argument of Proposition \ref{prop:II-2-2}, with module fractions in place
of ring fractions.
\end{proof}

\begin{proposition}[Quasi-coherent sheaves on affines]
\label{prop:II-5-2}
\source{hartshorne-algebraic-geometry:II.5.2}
\dcref{II.5.2}
If $X=\operatorname{Spec}A$, then $M\mapsto\widetilde M$ is an equivalence between $A$-modules and
quasi-coherent sheaves, with quasi-inverse $\Gamma(X,-)$. For $A\to B$,
$f^*\widetilde M=\widetilde{M\otimes_AB}$.
\uses{prop:II-5-1}
\end{proposition}
\begin{proof}
The unit and counit are isomorphisms on each D(f) by Proposition \ref{prop:II-5-1}; localization
commutes with tensor products, giving the pullback formula.
\end{proof}

\begin{lemma}[Localization criterion]
\label{lem:II-5-3}
\source{hartshorne-algebraic-geometry:II.5.3}
\dcref{II.5.3}
If $f\in A$ and $\mathcal F$ is quasi-coherent on $\operatorname{Spec}A$, then
$\Gamma(D(f),\mathcal F)=\Gamma(X,\mathcal F)_f$.
\end{lemma}
\begin{proof}
It holds for associated sheaves and hence locally; restriction and gluing preserve the asserted
localization.
\end{proof}

\begin{proposition}[Local characterization]
\label{prop:II-5-4}
\source{hartshorne-algebraic-geometry:II.5.4}
\dcref{II.5.4}
An $\mathcal O_X$-module is quasi-coherent iff it is locally the cokernel of a morphism between free
modules; on a Noetherian scheme it is coherent iff the free modules may be chosen finite rank.
\end{proposition}
\begin{proof}
For $\widetilde M$ choose a presentation by free $A$-modules and localize. Conversely, local cokernel
descriptions give localized sections on basic opens and hence associated sheaves.
\end{proof}

\begin{corollary}[Affine equivalence]
\label{cor:II-5-5}
\source{hartshorne-algebraic-geometry:II.5.5}
\dcref{II.5.5}
The functors $M\mapsto\widetilde M$ and $\Gamma$ are mutually inverse on quasi-coherent sheaves.
\uses{prop:II-5-1}
\end{corollary}
\begin{proof}
The unit $M\to\Gamma(\operatorname{Spec}A,\widetilde M)$ and counit
$\widetilde{\Gamma(X,\mathcal F)}\to\mathcal F$ are isomorphisms on every basic open by
Proposition~\ref{prop:II-5-1}; the basic opens form a basis, so the maps are isomorphisms.
\end{proof}

\begin{proposition}[Exactness on affines]
\label{prop:II-5-6}
\source{hartshorne-algebraic-geometry:II.5.6}
\dcref{II.5.6}
For an affine scheme, a sequence of quasi-coherent sheaves is exact iff the sequence of global-section
modules is exact.
\end{proposition}
\begin{proof}
One implication follows from stalkwise exactness. For the converse, localize the exact sequence of
modules at every prime.
\end{proof}

\begin{proposition}[Operations on quasi-coherent modules]
\label{prop:II-5-7}
\source{hartshorne-algebraic-geometry:II.5.7}
\dcref{II.5.7}
Kernels, cokernels, images, and extensions of quasi-coherent sheaves are quasi-coherent; on a
Noetherian scheme the same holds for coherent sheaves.
\uses{prop:II-5-6}
\end{proposition}
\begin{proof}
Reduce to an affine. Kernels, cokernels and images follow from the module equivalence. For an extension,
global sections are exact by Proposition \ref{prop:II-5-6}; sheafify the resulting module sequence and
use the five lemma. Finite generation is preserved over Noetherian rings.
\end{proof}

\begin{proposition}[Direct and inverse image of modules]
\label{prop:II-5-8}
\source{hartshorne-algebraic-geometry:II.5.8}
\dcref{II.5.8}
For $f:X\to Y$, inverse images of quasi-coherent sheaves are quasi-coherent. Direct images are
quasi-coherent when X is Noetherian or f is quasi-compact and separated; in the Noetherian case
coherence is preserved by $f^*$.
\uses{prop:II-5-2,prop:II-5-7}
\end{proposition}
\begin{proof}
The affine cases follow from Proposition \ref{prop:II-5-2}. For direct image, cover a quasi-compact
source by finitely many affines; separatedness makes pairwise intersections affine, while Noetherianity
makes them quasi-compact and finitely coverable by affines. The equalizer sequence for compatible
sections and Proposition \ref{prop:II-5-7} give quasi-coherence.
\end{proof}

\begin{proposition}[Ideal sheaves and closed subschemes]
\label{prop:II-5-9}
\source{hartshorne-algebraic-geometry:II.5.9}
\dcref{II.5.9}
The ideal sheaf of a closed subscheme is quasi-coherent, coherent when X is Noetherian, and every
quasi-coherent ideal sheaf is the ideal sheaf of a unique closed subscheme.
\uses{prop:II-5-8,prop:II-5-7}
\end{proposition}
\begin{proof}
The ideal is the kernel of $\mathcal O_X\to i_*\mathcal O_Y$, so use Proposition \ref{prop:II-5-8}
and Proposition \ref{prop:II-5-7}. Conversely, locally write the ideal as $\widetilde{\mathfrak a}$;
the quotient defines $\operatorname{Spec}(A/\mathfrak a)$, and the affine constructions glue.
\end{proof}

\begin{corollary}[Affine closed subschemes]
\label{cor:II-5-10}
\source{hartshorne-algebraic-geometry:II.5.10}
\dcref{II.5.10}
Ideals of $A$ are in bijection with closed subschemes of $\operatorname{Spec}A$; every such closed subscheme
is affine.
\uses{prop:II-5-9}
\end{corollary}
\begin{proof}
For an ideal $\mathfrak a\subset A$, the quotient map $A\to A/\mathfrak a$ gives the closed immersion
$\operatorname{Spec}(A/\mathfrak a)\hookrightarrow\operatorname{Spec}A$.  Conversely, a closed subscheme
has a quasi-coherent ideal sheaf by Proposition~\ref{prop:II-5-9}; on the affine scheme this is
$\widetilde{\mathfrak a}$ for the ideal of its global sections, and the quotient is
$\operatorname{Spec}(A/\mathfrak a)$.  The two constructions are inverse.
\end{proof}

\begin{proposition}[Proj modules]
\label{prop:II-5-11}
\source{hartshorne-algebraic-geometry:II.5.11}
\dcref{II.5.11}
For graded $S$ and $M$, the associated sheaf on $\operatorname{Proj}S$ satisfies
$\Gamma(D_+(f),\widetilde M)=M_{(f)}$. It is quasi-coherent.
\end{proposition}
\begin{proof}
Apply the affine construction to each degree-zero localization and glue on overlaps.
\end{proof}

\begin{proposition}[Twisting sheaves]
\label{prop:II-5-12}
\source{hartshorne-algebraic-geometry:II.5.12}
\dcref{II.5.12}
The sheaf $\mathcal O_X(n)=\widetilde{S(n)}$ on $X=\operatorname{Proj}S$ is invertible when $X$ is
covered by standard affine charts; $\mathcal O_X(m)\otimes\mathcal O_X(n)\simeq\mathcal O_X(m+n)$.
\end{proposition}
\begin{proof}
On $D_+(f)$, multiplication of homogeneous fractions gives the trivialization; the transition functions
are units and are compatible with tensor products.
\end{proof}

\begin{proposition}[Generation and projective embedding]
\label{prop:II-5-13}
\source{hartshorne-algebraic-geometry:II.5.13}
\dcref{II.5.13}
For $S=A[x_0,\ldots,x_r]$ and $X=\operatorname{Proj}S=\mathbb P^r_A$, $\mathcal O_X(1)$ is invertible and
its coordinate sections generate it. A coherent sheaf on a projective scheme is generated after a
sufficiently positive twist.
\end{proposition}
\begin{proof}
The coordinate sections generate on each $D_+(x_i)$. For a coherent sheaf, choose finite affine
generators on the standard cover, clear denominators after a common twist, and use the finite cover
to obtain global generators.
\end{proof}

\begin{lemma}[Twisted localization]
\label{lem:II-5-14}
\source{hartshorne-algebraic-geometry:II.5.14}
\dcref{II.5.14}
Let $\mathcal L$ be invertible on $X$, let $s\in\Gamma(X,\mathcal L)$, and let $X_s$ be the
open set where $s$ is nonvanishing. If $X$ is quasi-compact and a section of a quasi-coherent
$\mathcal F$ vanishes on $X_s$, then some power of $s$ kills it. If $X$ has a finite affine cover
on which $\mathcal L$ is free and whose pairwise intersections are quasi-compact, every section
of $\mathcal F|_{X_s}$ becomes the restriction of a global section of $\mathcal F\otimes\mathcal L^n$
after multiplying by a sufficiently high power of $s$.
\end{lemma}
\begin{proof}
On each affine trivializing open this is the localization criterion for modules; quasi-compactness
allows finitely many powers to be replaced by one common power.
\end{proof}

\begin{proposition}[Proj reconstruction]
\label{prop:II-5-15}
\source{hartshorne-algebraic-geometry:II.5.15}
\dcref{II.5.15}
If $S$ is generated by $S_1$ over $S_0$, $X=\operatorname{Proj}S$, and $\mathcal F$ is
quasi-coherent, then the natural map $\widetilde{\Gamma_*(\mathcal F)}\to\mathcal F$ is an isomorphism.
\end{proposition}
\begin{proof}
On each standard affine $D_+(f)$, apply the preceding lemma to identify sections with the degree-zero
localization of $\Gamma_*(\mathcal F)$; these identifications agree on overlaps.
\end{proof}

\begin{corollary}[Projective schemes as Proj]
\label{cor:II-5-16}
\source{hartshorne-algebraic-geometry:II.5.16}
\dcref{II.5.16}
If $Y$ is a closed subscheme of $\mathbb P^r_A$, it is $\operatorname{Proj}(S/I)$ for a homogeneous
ideal $I\subset S=A[x_0,\ldots,x_r]$. A scheme over $\operatorname{Spec}A$ is projective iff it is
isomorphic to $\operatorname{Proj}S$ for a graded $S$ with $S_0=A$ finitely generated by $S_1$.
\uses{prop:II-5-15}
\end{corollary}
\begin{proof}
The ideal sheaf of $Y$ is quasi-coherent; Proposition~\ref{prop:II-5-15} identifies it with the
sheaf associated to its homogeneous module of global twisted sections. The converse follows by
presenting $S$ as a quotient of a polynomial ring.
\end{proof}

\begin{theorem}[Serre]
\label{thm:II-5-17}
\source{hartshorne-algebraic-geometry:II.5.17}
\dcref{II.5.17}
If $X$ is projective over a Noetherian ring $A$ and $\mathcal F$ is coherent, then
$\mathcal F(n)$ is generated by global sections for $n\gg0$ and $H^i(X,\mathcal F(n))=0$ for $i>0,n\gg0$.
\end{theorem}
\begin{proof}
Resolve $\mathcal F$ by finite sums of $\mathcal O_X(-m)$ on projective space. The Cech complex of the
standard affine cover computes cohomology; multiplication by powers of the homogeneous coordinates
makes its positive-degree cohomology vanish in high twists. Induct on the length of the resolution.
\end{proof}

\begin{corollary}[Coherent sheaves are quotients of twists]
\label{cor:II-5-18}
\source{hartshorne-algebraic-geometry:II.5.18}
\dcref{II.5.18}
On a projective scheme over a Noetherian ring, every coherent sheaf is a quotient of a finite direct
sum of twisted structure sheaves $\mathcal O_X(n)$.
\end{corollary}
\begin{proof}
By Serre, a sufficiently positive twist is generated by finitely many global sections; the resulting
surjection from a finite sum of $\mathcal O_X$ becomes the assertion after twisting back.
\end{proof}

\begin{theorem}[Finiteness of global sections]
\label{thm:II-5-19}
\source{hartshorne-algebraic-geometry:II.5.19}
\dcref{II.5.19}
If $X$ is projective over a field $k$ and $\mathcal F$ is coherent, then $\Gamma(X,\mathcal F)$ is finite-dimensional
over k. If X is projective over a finitely generated k-algebra A, it is a finite A-module.
\end{theorem}
\begin{proof}
For a graded presentation, reduce to twists on $\operatorname{Proj}S$. The graded module formed from
all sufficiently twisted global sections is integral over the coordinate ring and hence finite by
finiteness of integral closure; each graded piece is finite. The affine pushforward statement follows
from the affine module criterion.
\end{proof}

\begin{corollary}[Coherent direct image]
\label{cor:II-5-20}
\source{hartshorne-algebraic-geometry:II.5.20}
\dcref{II.5.20}
If $f:X\to Y$ is projective of finite type over a field and $\mathcal F$ is coherent, then
$f_*\mathcal F$ is coherent.
\end{corollary}
\begin{proof}
The assertion is affine-local on $Y$ and follows from finite generation of global sections over each
finitely generated coordinate ring.
\end{proof}

\section{Divisors}

\begin{definition}[Weil divisors]
\label{def:II-6-weil}
\source{hartshorne-algebraic-geometry:II.6}
On a Noetherian integral scheme regular in codimension one, a prime divisor is an integral closed
subscheme of codimension one. A Weil divisor is a finite sum $D=\sum n_Y Y$. The local ring at the
generic point of $Y$ is a DVR, giving $v_Y:K(X)^*\to\mathbb Z$.
\end{definition}

\begin{proposition}[Principal divisors]
\label{prop:II-6-1}
\source{hartshorne-algebraic-geometry:II.6.1}
\dcref{II.6.1}
For $f\in K(X)^*$, $(f)=\sum_Yv_Y(f)Y$ has finite support. Principal divisors form a subgroup of Weil
divisors.
\end{proposition}
\begin{proof}
On an affine where $f$ is regular, positive valuations occur only at components of $V(f)$; apply the same
argument to $f^{-1}$ on an affine where it is regular. The valuation identity
$v_Y(fg)=v_Y(f)+v_Y(g)$ proves additivity.
\end{proof}

\begin{definition}[Linear equivalence and class group]
\label{def:II-6-class}
\source{hartshorne-algebraic-geometry:II.6}
$D\sim D'$ when $D-D'=(f)$. The class group is $\operatorname{Cl}X=\operatorname{Div}X/\operatorname{Prin}X$.
\end{definition}

\begin{theorem}[Factoriality and the divisor class group]
\label{thm:II-6-2}
\source{hartshorne-algebraic-geometry:II.6.2}
\dcref{II.6.2}
Let $A$ be a noetherian domain and $X=\operatorname{Spec}A$.  Then $A$ is a
unique factorization domain if and only if $X$ is normal and $\operatorname{Cl}X=0$.
\end{theorem}
\begin{proof}
A UFD is integrally closed, hence $X$ is normal.  If every height-one prime is principal, every
prime divisor is principal and $\operatorname{Cl}X=0$.  Conversely, if $X$ is normal and its class
group is zero, a height-one prime divisor is linearly equivalent to zero, hence is the divisor of
an element of $K(X)^*$; valuation one at that prime and nonnegative valuations elsewhere show that
the element lies in $A$ and generates the prime.  The height-one criterion for a noetherian domain
then gives factoriality.
\end{proof}

\begin{proposition}[Intersection of height-one localizations]
\label{prop:II-6-3A}
\source{hartshorne-algebraic-geometry:II.6.3A}
\dcref{II.6.3A}
If $A$ is an integrally closed noetherian domain, then
\[
A=\bigcap_{\operatorname{ht}\mathfrak p=1}A_{\mathfrak p}
\]
inside the fraction field of $A$, where the intersection runs over all height-one prime ideals.
\end{proposition}
\begin{proof}
This is the standard intersection theorem for Krull domains; Hartshorne cites Matsumura for the
result and does not prove it in the text.
\end{proof}

\begin{definition}[Cartier divisors]
\label{def:II-6-cartier}
\source{hartshorne-algebraic-geometry:II.6}
A Cartier divisor is a compatible collection of rational functions modulo units. It determines an
invertible sheaf $\mathcal O_X(D)$; addition corresponds to tensor product and principal divisors give
$\mathcal O_X((f))\simeq\mathcal O_X$.
\end{definition}

\begin{proposition}[Divisors on projective space]
\label{prop:II-6-4}
\source{hartshorne-algebraic-geometry:II.6.4}
\dcref{II.6.4}
For $n\geq1$, every divisor on $\mathbb P^n$ is linearly equivalent to a unique multiple of a
hyperplane; hence $\operatorname{Cl}(\mathbb P^n)\simeq\mathbb Z$ generated by a hyperplane.
\end{proposition}
\begin{proof}
On the standard affine cover, clearing denominators shows that a divisor of degree zero is principal;
the degree map sends a hyperplane to $1$ and detects the unique multiple.
\end{proof}

\begin{proposition}[Restriction of divisors]
\label{prop:II-6-5}
\source{hartshorne-algebraic-geometry:II.6.5}
\dcref{II.6.5}
If $Y$ is a proper closed subset of a noetherian integral scheme $X$, restriction induces an exact
sequence $\operatorname{Div}_Y(X)\to\operatorname{Cl}(X)\to\operatorname{Cl}(X\setminus Y)\to0$.
\end{proposition}
\begin{proof}
A divisor on the open complement closes to a divisor on $X$, and changing the closure changes it by
a divisor supported on $Y$; principal divisors restrict to principal divisors.
\end{proof}

\begin{proposition}[Affine-line invariance of the class group]
\label{prop:II-6-6}
\source{hartshorne-algebraic-geometry:II.6.6}
\dcref{II.6.6}
If $X$ satisfies the standing hypotheses of this section, then $X\times\mathbb A^1$ also satisfies
them, and
\[
\operatorname{Cl}X\xrightarrow{\sim}\operatorname{Cl}(X\times\mathbb A^1)
\]
via pullback along the projection.
\end{proposition}
\begin{proof}
The product is noetherian, integral, and separated.  Codimension-one points are either pullbacks of
codimension-one points of $X$ or dominate $X$; the latter have local rings which are localizations
of $K(X)[t]$ and hence principal divisor class.  Sending a divisor to the sum of the pullbacks of
its components defines the inverse on class groups; valuation computations show that both composites
are the identity.
\end{proof}

\begin{proposition}[Divisors on nonsingular curves]
\label{prop:II-6-7}
\source{hartshorne-algebraic-geometry:II.6.7}
\dcref{II.6.7}
For a complete nonsingular curve, every Weil divisor is Cartier, and the divisor of a nonzero
rational function is the sum of its valuations at the closed points with finite support.
\end{proposition}
\begin{proof}
The local rings at closed points are DVRs, so local equations exist and glue as Cartier data; finite
support follows by applying regularity to the function and its inverse on affine neighborhoods.
\end{proof}

\begin{proposition}[Finite maps of curves]
\label{prop:II-6-8}
\source{hartshorne-algebraic-geometry:II.6.8}
\dcref{II.6.8}
A nonconstant map of complete nonsingular curves is finite. For $Q\in Y$,
$f^*Q=\sum_{f(P)=Q}v_P(t)P$, where $t$ is a uniformizer at $Q$, and this is independent of $t$.
\end{proposition}
\begin{proof}
A nonconstant map has zero-dimensional fibres; properness of complete curves makes it finite. A changed
uniformizer differs by a unit, whose pullback has valuation zero. Integral closure over the DVR at Q
gives the displayed ramification multiplicities.
\end{proof}

\begin{proposition}[Degree pullback]
\label{prop:II-6-9}
\source{hartshorne-algebraic-geometry:II.6.9}
\dcref{II.6.9}
For a finite map of nonsingular complete curves, $\deg f^*D=(\deg f)\deg D$.
\end{proposition}
\begin{proof}
It suffices to use one point $Q$. The integral closure of $\mathcal O_{Y,Q}$ is a free module of rank
[K(X):K(Y)]; its quotient by a uniformizer has length equal to that rank and also equals the sum
of the ramification indices.
\end{proof}

\begin{corollary}[Principal divisors have degree zero]
\label{cor:II-6-10}
\source{hartshorne-algebraic-geometry:II.6.10}
\dcref{II.6.10}
A principal divisor on a complete nonsingular curve has degree zero.
\end{corollary}
\begin{proof}
Apply the valuation theorem to multiplication by f on the finite-dimensional Riemann--Roch function
space, or compare zeros and poles through a finite map to $\mathbb P^1$; the total valuation sum is zero.
\end{proof}

\begin{proposition}[Cartier and Weil divisors]
\label{prop:II-6-11}
\source{hartshorne-algebraic-geometry:II.6.11}
\dcref{II.6.11}
On an integral separated noetherian normal scheme, Cartier divisors map to Weil divisors by taking
the valuations at codimension-one points, and this map identifies principal Cartier divisors with
principal Weil divisors.
\end{proposition}
\begin{proof}
At a codimension-one point the local ring is a DVR, so a local Cartier equation has a well-defined
order. Compatibility on overlaps gives a Weil divisor; normality makes the construction injective.
\end{proof}

\begin{proposition}[Invertible sheaves and duals]
\label{prop:II-6-12}
\source{hartshorne-algebraic-geometry:II.6.12}
\dcref{II.6.12}
For an invertible sheaf $\mathcal L$, its dual $\mathcal L^{-1}$ is invertible and
$\mathcal L\otimes\mathcal L^{-1}\simeq\mathcal O_X$; isomorphism classes form the Picard group
$\operatorname{Pic}(X)$ under tensor product.
\end{proposition}
\begin{proof}
Trivialize $\mathcal L$ on an open cover. The dual is locally free of rank one, and evaluation gives
the inverse; tensor product gives the group law.
\end{proof}

\begin{proposition}[Cartier divisors and the Picard group]
\label{prop:II-6-13}
\source{hartshorne-algebraic-geometry:II.6.13}
\dcref{II.6.13}
The map $D\mapsto\mathcal O_X(D)$ induces an isomorphism from Cartier divisors modulo principal
divisors to $\operatorname{Pic}(X)$.
\end{proposition}
\begin{proof}
Local equations give the invertible sheaf and changing equations by units gives an isomorphic sheaf.
Conversely, local generators of an invertible sheaf have rational ratios defining inverse Cartier data.
\end{proof}

\begin{corollary}[Cartier-to-Weil injectivity]
\label{cor:II-6-14}
\source{hartshorne-algebraic-geometry:II.6.14}
\dcref{II.6.14}
On an integral scheme, the map from Cartier divisors modulo principal divisors to the divisor class
group is injective.
\end{corollary}
\begin{proof}
If the associated Weil divisor is principal, divide the Cartier data by that rational function; the
resulting local equations are units, so the Cartier divisor is principal.
\end{proof}

\begin{proposition}[Effective Cartier divisors]
\label{prop:II-6-15}
\source{hartshorne-algebraic-geometry:II.6.15}
\dcref{II.6.15}
An effective Cartier divisor is locally defined by a non-zero-divisor. Its ideal sheaf is invertible,
and the quotient defines a closed subscheme of pure codimension one.
\end{proposition}
\begin{proof}
Multiplication by each local equation is injective and identifies the ideal with a free rank-one
module; the local quotients glue to the asserted closed subscheme.
\end{proof}

\begin{corollary}[Locally factorial schemes]
\label{cor:II-6-16}
\source{hartshorne-algebraic-geometry:II.6.16}
\dcref{II.6.16}
On a locally factorial scheme every Weil divisor is Cartier, and the Cartier and Weil class groups
are canonically isomorphic.
\end{corollary}
\begin{proof}
Codimension-one equations extend on factorial affine neighborhoods; apply the Cartier--Weil
identification above.
\end{proof}

\begin{corollary}[Class group of projective space]
\label{cor:II-6-17}
\source{hartshorne-algebraic-geometry:II.6.17}
\dcref{II.6.17}
For $n\geq1$, $\operatorname{Pic}(\mathbb P^n)\simeq\mathbb Z$, generated by $\mathcal O(1)$.
\uses{prop:II-6-4}
\end{corollary}
\begin{proof}
Combine Proposition~\ref{prop:II-6-4} with the Cartier-to-Weil identification.
\end{proof}

\begin{proposition}[Effective Cartier divisors and closed subschemes]
\label{prop:II-6-18}
\source{hartshorne-algebraic-geometry:II.6.18}
\dcref{II.6.18}
Effective Cartier divisors on $X$ are in bijection with closed subschemes whose ideal sheaf is
invertible and locally generated by a non-zero-divisor.
\end{proposition}
\begin{proof}
The two constructions are $D\mapsto\mathcal O_X(-D)$ and an invertible ideal with a local
non-zero-divisor generator to its vanishing divisor; the constructions are inverse.
\end{proof}

\section{Projective Morphisms}

\begin{definition}[Projective morphism]
\label{def:II-7-projective}
\source{hartshorne-algebraic-geometry:II.7}
A morphism is projective if it factors as a closed immersion $X\hookrightarrow\mathbb P^n_Y$. An
invertible sheaf is relatively very ample when it is the pullback of $\mathcal O(1)$ under such an embedding.
\end{definition}

\begin{theorem}[Morphisms to projective space]
\label{thm:II-7-1}
\source{hartshorne-algebraic-geometry:II.7.1}
\dcref{II.7.1}
A morphisms $f:X\to\mathbb P^n_Y$ is equivalent to an invertible sheaf $\mathcal L$ on $X$
together with generating global sections $s_0,\ldots,s_n$, up to the natural change of basis.
\end{theorem}
\begin{proof}
Pull back $\mathcal O(1)$ and its coordinate sections. Conversely, the ratios of the generating
sections define compatible maps to the standard affine charts and glue to $f$.
\end{proof}

\begin{proposition}[Closed immersion criterion]
\label{prop:II-7-2}
\source{hartshorne-algebraic-geometry:II.7.2}
\dcref{II.7.2}
A morphism defined by generating sections of an invertible sheaf is a closed immersion precisely
when the induced maps on the standard affine opens identify their coordinate rings with the relevant
affine coordinate rings.
\end{proposition}
\begin{proof}
Check the assertion on the standard affine cover of projective space; a morphism is a closed immersion
if and only if it is so on an affine open cover of the target.
\end{proof}

\begin{proposition}[Projective embedding criterion]
\label{prop:II-7-3}
\source{hartshorne-algebraic-geometry:II.7.3}
\dcref{II.7.3}
For a finite-type scheme over a field, a finite-dimensional space of sections of an invertible sheaf
defines a closed immersion into projective space iff it separates points and tangent vectors.
\end{proposition}
\begin{proof}
Point separation makes the map injective and an immersion on residue fields; tangent-vector separation
makes the induced local maps surjective on cotangent spaces. The affine criterion then gives a closed
immersion.
\end{proof}

\begin{lemma}[Finite generation from local maps]
\label{lem:II-7-4}
\source{hartshorne-algebraic-geometry:II.7.4}
\dcref{II.7.4}
If a local homomorphism of noetherian local rings induces a finite extension on residue fields and
the target is generated by finitely many elements over the image, then the corresponding morphism is
finite on a neighborhood.
\end{lemma}
\begin{proof}
Choose finitely many generators and clear denominators in the localizations; Nakayama's lemma gives
finite generation after shrinking.
\end{proof}

\begin{proposition}[Ample invertible sheaves]
\label{prop:II-7-5}
\source{hartshorne-algebraic-geometry:II.7.5}
\dcref{II.7.5}
For an invertible sheaf $\mathcal L$ on a noetherian scheme, the following are equivalent: $\mathcal L$
is ample; every coherent sheaf becomes generated by global sections after a positive twist; and a
sufficiently high power of $\mathcal L$ gives a projective embedding on every affine neighborhood.
\end{proposition}
\begin{proof}
The implications follow by applying the finite-generation lemma to affine standard opens and using
the point-and-tangent separation criterion.
\end{proof}

\begin{theorem}[Projective implies proper]
\label{thm:II-7-6}
\source{hartshorne-algebraic-geometry:II.7.6}
\dcref{II.7.6}
Every projective morphism is proper.
\end{theorem}
\begin{proof}
The projection $\mathbb P^n_Y\to Y$ satisfies the valuative criterion by scaling homogeneous coordinates
so one is a unit in the valuation ring. Closed immersions and compositions preserve properness.
\end{proof}

\begin{proposition}[Divisors of zeros and linear systems]
\label{prop:II-7-7}
\source{hartshorne-algebraic-geometry:II.7.7}
\dcref{II.7.7}
A nonzero section of an invertible sheaf determines an effective Cartier divisor of zeros. The
complete linear system of an invertible sheaf is the projective space of its nonzero global sections,
and its associated map is the map defined by those sections.
\end{proposition}
\begin{proof}
Locally trivialize the sheaf; the section is a non-zero-divisor equation and its local equations differ
by units. Ratios of a basis of global sections give the associated projective map.
\end{proof}

\begin{lemma}[Base-point-free linear systems]
\label{lem:II-7-8}
\source{hartshorne-algebraic-geometry:II.7.8}
\dcref{II.7.8}
A linear system has no base points precisely when its sections generate the associated invertible
sheaf; in that case it defines a morphism to projective space.
\uses{thm:II-7-1}
\end{lemma}
\begin{proof}
The stalkwise generation condition is exactly the absence of a common zero, and Theorem~\ref{thm:II-7-1}
then constructs the morphism.
\end{proof}

\begin{lemma}[Proj invariant under twisting]
\label{lem:II-7-9}
\source{hartshorne-algebraic-geometry:II.7.9}
\dcref{II.7.9}
For a graded ring $S$, $\operatorname{Proj}S\simeq\operatorname{Proj}S(d)$ after replacing the grading
by a positive twist; the standard opens and their structure sheaves are unchanged.
\end{lemma}
\begin{proof}
The homogeneous prime ideals and degree-zero localizations are unchanged by a common shift of degrees.
\end{proof}

\begin{theorem}[Relative Proj is projective]
\label{thm:II-7-10}
\source{hartshorne-algebraic-geometry:II.7.10}
\dcref{II.7.10}
If $S$ is a finitely generated graded algebra over $A=S_0$, then $\operatorname{Proj}S\to\operatorname{Spec}A$
is projective.
\end{theorem}
\begin{proof}
Choose homogeneous generators in degree one after a twist; the resulting graded surjection from a
polynomial ring gives a closed immersion into projective space over $A$.
\end{proof}

\begin{proposition}[Relative Proj of a symmetric algebra]
\label{prop:II-7-11}
\source{hartshorne-algebraic-geometry:II.7.11}
\dcref{II.7.11}
For a quasi-coherent module $\mathcal E$ on $Y$, $\operatorname{Proj}\operatorname{Sym}(\mathcal E)$
represents invertible quotients of the pullback of $\mathcal E$.
\end{proposition}
\begin{proof}
This is the affine universal property of Proj applied on a cover of $Y$, with the local constructions
glued by functoriality.
\end{proof}

\begin{proposition}[Universal property of relative Proj]
\label{prop:II-7-12}
\source{hartshorne-algebraic-geometry:II.7.12}
\dcref{II.7.12}
Morphisms $T\to\operatorname{Proj}\operatorname{Sym}(\mathcal E)$ over $Y$ are equivalent to pairs
consisting of an invertible sheaf on $T$ and a surjection from the pullback of $\mathcal E$.
\end{proposition}
\begin{proof}
The universal quotient gives the forward map; local generators of a quotient give homogeneous
coordinates and glue to the inverse.
\end{proof}

\begin{proposition}[Blow-up ideal]
\label{prop:II-7-13}
\source{hartshorne-algebraic-geometry:II.7.13}
\dcref{II.7.13}
For the blow-up $\operatorname{Bl}_\mathcal I X=\operatorname{Proj}\bigoplus_{n\geq0}\mathcal I^n$,
the inverse image of $\mathcal I$ is an invertible ideal sheaf.
\end{proposition}
\begin{proof}
On a standard Proj chart one generator of $\mathcal I$ becomes a unit times the tautological degree-one
generator, so the inverse image is locally principal and hence invertible.
\end{proof}

\begin{proposition}[Universal property of blow-up]
\label{prop:II-7-14}
\source{hartshorne-algebraic-geometry:II.7.14}
\dcref{II.7.14}
The blow-up of $X$ along $\mathcal I$ is universal among morphisms $g:Y\to X$ for which
$g^{-1}\mathcal I\mathcal O_Y$ is invertible.
\end{proposition}
\begin{proof}
An invertible pullback ideal gives a graded algebra map from the Rees algebra to the section algebra
of its powers, hence a unique map to Proj; the tautological ideal proves uniqueness.
\end{proof}

\begin{corollary}[Functoriality of blow-up]
\label{cor:II-7-15}
\source{hartshorne-algebraic-geometry:II.7.15}
\dcref{II.7.15}
If a morphism pulls one ideal sheaf into another, it induces the corresponding morphism of blow-ups;
in particular blow-up commutes with base change in the stated ideal-theoretic sense.
\end{corollary}
\begin{proof}
Apply the universal property to the pullback of the tautological invertible ideal.
\end{proof}

\begin{proposition}[Blow-ups of varieties]
\label{prop:II-7-16}
\source{hartshorne-algebraic-geometry:II.7.16}
\dcref{II.7.16}
The blow-up of a variety along a coherent ideal is projective over the variety and is again a variety.
\end{proposition}
\begin{proof}
Finite generation of the Rees algebra gives projectivity by relative Proj. On affine charts the Rees
algebra is a domain, and the resulting scheme is integral and separated.
\end{proof}

\begin{theorem}[Factorization through a blow-up]
\label{thm:II-7-17}
\source{hartshorne-algebraic-geometry:II.7.17}
\dcref{II.7.17}
Every projective morphism of noetherian schemes factors as a blow-up followed by a closed immersion
into projective space (after choosing a suitable coherent ideal).
\end{theorem}
\begin{proof}
Choose a relatively ample invertible sheaf and finitely many generating sections. The associated
graded algebra gives a Rees algebra after clearing denominators; the universal property of the blow-up
then supplies the factorization.
\end{proof}

\section{Differentials}

\begin{definition}[Kahler differentials]
\label{def:II-8-differentials}
\source{hartshorne-algebraic-geometry:II.8}
For $A\to B$, $\Omega_{B/A}$ represents $A$-derivations: it is generated by symbols $db$ with
$d(b+b')=db+db'$, $d(bb')=b\,db'+b'\,db$, and $da=0$.
\end{definition}

\begin{proposition}[Diagonal description]
\label{prop:II-8-1A}
\source{hartshorne-algebraic-geometry:II.8.1A}
\dcref{II.8.1A}
Let $B$ be an $A$-algebra, let $I$ be the kernel of the multiplication map
$B\otimes_A B\to B$, and define $db=1\otimes b-b\otimes1$ modulo $I^2$.  Then
$(I/I^2,d)$ is a module of relative differentials for $B/A$.
\end{proposition}
\begin{proof}
The displayed map satisfies the derivation identities.  Any $A$-derivation $B\to M$ induces a
$B$-linear map $I/I^2\to M$ by $b\otimes b'\mapsto b\,db'$, and this is the unique such map.
Hartshorne cites Matsumura for the construction.
\end{proof}

\begin{proposition}[Base change and localization]
\label{prop:II-8-2A}
\source{hartshorne-algebraic-geometry:II.8.2A}
\dcref{II.8.2A}
If $A'$ and $B$ are $A$-algebras and $B'=B\otimes_AA'$, then
$\Omega_{B'/A'}\simeq\Omega_{B/A}\otimes_BB'$.  If $S$ is a multiplicative system in $B$, then
$\Omega_{S^{-1}B/A}\simeq S^{-1}\Omega_{B/A}$.
\end{proposition}
\begin{proof}
Extend derivations uniquely across the base change and across inverses; the two universal properties
give the displayed isomorphisms.  Hartshorne cites Matsumura for the result.
\end{proof}

\begin{proposition}[First exact sequence]
\label{prop:II-8-3A}
\source{hartshorne-algebraic-geometry:II.8.3A}
\dcref{II.8.3A}
For ring maps $A\to B\to C$ there is a natural exact sequence of $C$-modules
\[
\Omega_{B/A}\otimes_BC\longrightarrow\Omega_{C/A}\longrightarrow\Omega_{C/B}\longrightarrow0.
\]
\end{proposition}
\begin{proof}
The first map is induced by composition of derivations.  A $B$-derivation is exactly an $A$-derivation
which kills the image of $B$, so the universal property gives exactness; the source cites Matsumura.
\end{proof}

\begin{proposition}[Second exact sequence]
\label{prop:II-8-4A}
\source{hartshorne-algebraic-geometry:II.8.4A}
\dcref{II.8.4A}
If $B$ is an $A$-algebra, $I\subset B$ an ideal, and $C=B/I$, then there is an exact sequence
\[
I/I^2\xrightarrow{\delta}\Omega_{B/A}\otimes_BC\longrightarrow\Omega_{C/A}\longrightarrow0,
\]
where $\delta(\bar b)=db\otimes1$.
\end{proposition}
\begin{proof}
Apply the diagonal description to the quotient map and use the universal property of $\Omega$;
the kernel consists precisely of differentials of elements of $I$, modulo $I^2$.  Hartshorne cites
Matsumura for this exact sequence.
\end{proof}

\begin{corollary}[Finite generation of differentials]
\label{cor:II-8-5}
\source{hartshorne-algebraic-geometry:II.8.5}
\dcref{II.8.5}
If $B$ is a finitely generated $A$-algebra, or a localization of one, then $\Omega_{B/A}$ is a
finitely generated $B$-module.
\end{corollary}
\begin{proof}
Present $B$ as a quotient of a polynomial algebra.  Its module of differentials is generated by the
coordinate differentials, and the second exact sequence and localization preserve finite generation.
\end{proof}

\begin{theorem}[Separability criterion]
\label{thm:II-8-6A}
\source{hartshorne-algebraic-geometry:II.8.6A}
\dcref{II.8.6A}
For a finitely generated extension field $K/k$, $\dim_K\Omega_{K/k}\ge\operatorname{trdeg}_kK$, with
equality if and only if $K/k$ is separably generated.
\end{theorem}
\begin{proof}
Choose a transcendence basis and apply the exact sequence to the finite algebraic extension over the
rational function field.  The basis differentials are independent, and the algebraic part contributes
no additional dimension exactly when it is separable; the source cites Matsumura.
\end{proof}

\begin{proposition}[Residue-field differentials]
\label{prop:II-8-7}
\source{hartshorne-algebraic-geometry:II.8.7}
\dcref{II.8.7}
Let $B$ be a local ring containing a field $k$ isomorphic to its residue field $B/\mathfrak m$.  Then
$\delta:\mathfrak m/\mathfrak m^2\to\Omega_{B/k}\otimes_Bk$ is an isomorphism.
\end{proposition}
\begin{proof}
The second exact sequence identifies the cokernel with $\Omega_{k/k}=0$.  A $k$-derivation on $B$
whose restriction to $\mathfrak m$ vanishes modulo $\mathfrak m^2$ is zero, proving injectivity.
\end{proof}

\begin{theorem}[Regularity via differentials]
\label{thm:II-8-8}
\source{hartshorne-algebraic-geometry:II.8.8}
\dcref{II.8.8}
Let $B$ be a local ring containing a perfect field $k$ isomorphic to its residue field, and assume
$B$ is a localization of a finitely generated $k$-algebra.  Then $\Omega_{B/k}$ is free of rank
$\dim B$ if and only if $B$ is a regular local ring.
\uses{prop:II-8-7}
\end{theorem}
\begin{proof}
By Proposition~\ref{prop:II-8-7}, $\mathfrak m/\mathfrak m^2$ is the cotangent space.  The exact
sequence and separability criterion identify its dimension with the rank of $\Omega_{B/k}$; equality
with $\dim B$ is the regular-local criterion.
\end{proof}

\begin{lemma}[Free-module criterion]
\label{lem:II-8-9}
\source{hartshorne-algebraic-geometry:II.8.9}
\dcref{II.8.9}
If $A$ is a noetherian local integral domain with residue field $k$ and quotient field $K$, and $M$ is
a finite $A$-module satisfying $\dim_k(M\otimes_Ak)=\dim_K(M\otimes_AK)=r$, then $M$ is free of rank $r$.
\end{lemma}
\begin{proof}
Nakayama gives a surjection $A^r\to M$.  Tensoring with $K$ shows that its kernel has rank zero;
torsion-freeness forces the kernel to vanish, so the map is an isomorphism.
\end{proof}

\begin{definition}[Sheaf of relative differentials]
\label{def:II-8-sheaf-differentials}
\source{hartshorne-algebraic-geometry:II.8}
For a morphism $f:X\to Y$ of schemes, the sheaf of relative differentials $\Omega_{X/Y}$ is obtained
on affine opens from the modules $\Omega_{B/A}$ and glued using localization; equivalently, it is
the pullback of the conormal sheaf of the diagonal $X\to X\times_YX$.
\end{definition}

\begin{proposition}[Base change for sheaf differentials]
\label{prop:II-8-10}
\source{hartshorne-algebraic-geometry:II.8.10}
\dcref{II.8.10}
For $f:X\to Y$ and $g:Y'\to Y$, with $X'=X\times_YY'$, one has
$\Omega_{X'/Y'}\simeq g'^*\Omega_{X/Y}$, where $g':X'\to X$ is the projection.
\uses{prop:II-8-2A}
\end{proposition}
\begin{proof}
Work on affine charts and apply Proposition~\ref{prop:II-8-2A}; the local isomorphisms agree on
overlaps by the universal property and therefore glue.
\end{proof}

\begin{proposition}[Transitivity sequence]
\label{prop:II-8-11}
\source{hartshorne-algebraic-geometry:II.8.11}
\dcref{II.8.11}
For morphisms $X\xrightarrow{f}Y\xrightarrow{g}Z$, there is an exact sequence
\[
f^*\Omega_{Y/Z}\longrightarrow\Omega_{X/Z}\longrightarrow\Omega_{X/Y}\longrightarrow0.
\]
\end{proposition}
\begin{proof}
The assertion is the sheafification of the first exact sequence on affine charts.
\end{proof}

\begin{proposition}[Closed-subscheme sequence]
\label{prop:II-8-12}
\source{hartshorne-algebraic-geometry:II.8.12}
\dcref{II.8.12}
If $Z\hookrightarrow X$ is a closed subscheme with ideal sheaf $\mathcal I$ and $X\to Y$ is a
morphism, then there is an exact sequence on $Z$
\[
\mathcal I/\mathcal I^2\xrightarrow{\delta}\Omega_{X/Y}\otimes\mathcal O_Z\longrightarrow\Omega_{Z/Y}\longrightarrow0.
\]
\end{proposition}
\begin{proof}
Apply the second exact sequence on affine charts and glue.  The map sends the class of a local
equation to its differential restricted to $Z$.
\end{proof}

\begin{theorem}[Euler sequence]
\label{thm:II-8-13}
\source{hartshorne-algebraic-geometry:II.8.13}
\dcref{II.8.13}
On $\mathbb P^n_A$,
\[
0\to\Omega_{\mathbb P^n_A/A}\to\mathcal O(-1)^{\oplus(n+1)}
\xrightarrow{(x_0,\ldots,x_n)}\mathcal O\to0.
\]
\end{theorem}
\begin{proof}
On each standard chart, identify differentials of $x_j/x_i$ and check that the kernel of the Euler
map is generated by these relations. The local descriptions agree on overlaps and exactness is
stalkwise.
\end{proof}

\begin{theorem}[Localization of regular local rings]
\label{thm:II-8-14A}
\source{hartshorne-algebraic-geometry:II.8.14A}
\dcref{II.8.14A}
If $(A,\mathfrak m)$ is a regular local ring, every localization $A_\mathfrak p$ is regular.
\end{theorem}
\begin{proof}
A regular system of parameters remains a minimal generating set after localization; the
dimension and embedding-dimension inequality then give regularity at each prime.
\end{proof}

\begin{theorem}[Differential criterion for nonsingularity]
\label{thm:II-8-15}
\source{hartshorne-algebraic-geometry:II.8.15}
\dcref{II.8.15}
For a variety over a perfect field, a point is nonsingular if and only if the stalk of
$\Omega_{X/k}$ is free of rank $\dim X$.
\uses{thm:II-8-6A}
\end{theorem}
\begin{proof}
Use the conormal sequence and Theorem~\ref{thm:II-8-6A}; the cotangent-space dimension equals
the local embedding dimension, which equals the dimension precisely for a regular local ring.
\end{proof}

\begin{corollary}[Open dense nonsingular locus]
\label{cor:II-8-16}
\source{hartshorne-algebraic-geometry:II.8.16}
\dcref{II.8.16}
For a variety over a perfect field, the nonsingular locus is open and dense.
\end{corollary}
\begin{proof}
The rank condition on a finite presentation of $\Omega_{X/k}$ is open, and the generic point is
nonsingular by the differential criterion.
\end{proof}

\begin{theorem}[Conormal characterization]
\label{thm:II-8-17}
\source{hartshorne-algebraic-geometry:II.8.17}
\dcref{II.8.17}
A finite-type variety $X$ over a perfect field is nonsingular if and only if $\Omega_{X/k}$ is
locally free of rank $\dim X$; for a closed immersion $X\hookrightarrow Y$ between nonsingular
varieties, the conormal sequence is exact on the left as well.
\uses{thm:II-8-15}
\end{theorem}
\begin{proof}
The first assertion is Theorem~\ref{thm:II-8-15}. In local polynomial coordinates, regularity
makes the defining equations a regular sequence, so the conormal map is injective.
\end{proof}

\begin{theorem}[Bertini for hyperplane sections]
\label{thm:II-8-18}
\source{hartshorne-algebraic-geometry:II.8.18}
\dcref{II.8.18}
If $X\subset\mathbb P^n$ is nonsingular over an infinite perfect field, a general hyperplane
section of $X$ is nonsingular (or empty).
\end{theorem}
\begin{proof}
The hyperplanes tangent to $X$ form a proper closed subset of the dual projective space by the
dimension estimate. Outside it, the conormal criterion gives a nonsingular section.
\end{proof}

\begin{theorem}[Birational invariance of geometric genus]
\label{thm:II-8-19}
\source{hartshorne-algebraic-geometry:II.8.19}
\dcref{II.8.19}
The geometric genus of a nonsingular projective variety is invariant under birational equivalence.
\end{theorem}
\begin{proof}
Resolve the birational correspondence on a common open set; regular top differential forms extend
across codimension at least two on a nonsingular variety, so the spaces and dimensions agree.
\end{proof}

\begin{corollary}[Canonical sheaf of projective space]
\label{cor:II-8-20}
\source{hartshorne-algebraic-geometry:II.8.20}
\dcref{II.8.20}
$\bigwedge^n\Omega_{\mathbb P^n_A/A}\simeq\mathcal O_{\mathbb P^n_A}(-n-1)$.
\end{corollary}
\begin{proof}
Take the top exterior power of the Euler sequence; the determinant of
$\mathcal O(-1)^{\oplus(n+1)}$ is $\mathcal O(-n-1)$.
\end{proof}

\begin{theorem}[Cohen--Macaulay criteria]
\label{thm:II-8-21A}
\source{hartshorne-algebraic-geometry:II.8.21A}
\dcref{II.8.21A}
For a noetherian local ring, a quotient of a regular local ring by a regular sequence is
Cohen--Macaulay; equivalently, its depth equals its dimension.
\end{theorem}
\begin{proof}
The Koszul complex of a regular sequence is a finite free resolution; its depth computation gives
the equality, and localization preserves the property.
\end{proof}

\begin{theorem}[Serre normality criterion]
\label{thm:II-8-22A}
\source{hartshorne-algebraic-geometry:II.8.22A}
\dcref{II.8.22A}
A noetherian ring is normal if and only if it satisfies $(R_1)$, regularity at codimension-one
primes, and $(S_2)$, depth at least the minimum of $2$ and the local dimension.
\end{theorem}
\begin{proof}
Apply the intersection description of integral closure in height-one valuation rings and the depth
criterion for extension across codimension at least two.
\end{proof}

\begin{proposition}[Local complete intersections]
\label{prop:II-8-23}
\source{hartshorne-algebraic-geometry:II.8.23}
\dcref{II.8.23}
A local complete intersection in a nonsingular variety is Cohen--Macaulay and is regular in
codimension one.
\uses{thm:II-8-21A}
\end{proposition}
\begin{proof}
Its ideal is locally generated by a regular sequence, so Theorem~\ref{thm:II-8-21A} gives
Cohen--Macaulayness; the codimension-one localizations are hypersurfaces in regular local rings.
\end{proof}

\begin{theorem}[Blow-up along a local complete intersection]
\label{thm:II-8-24}
\source{hartshorne-algebraic-geometry:II.8.24}
\dcref{II.8.24}
The blow-up of a nonsingular variety along a local complete intersection is nonsingular.
\end{theorem}
\begin{proof}
On an affine chart the center is generated by a regular sequence; the blow-up charts are polynomial
extensions of the quotient by all but one generator, hence regular by the Jacobian criterion.
\end{proof}

\begin{theorem}[Separable scalar extension]
\label{thm:II-8-25A}
\source{hartshorne-algebraic-geometry:II.8.25A}
\dcref{II.8.25A}
If $K/k$ is a separable field extension and $X$ is a nonsingular variety over $k$, then
$X_K=X\times_k\operatorname{Spec}K$ is nonsingular.
\end{theorem}
\begin{proof}
The exact differential sequence for $k\subset K$ and separability give $\Omega_{K/k}=0$; base
change therefore preserves the locally free differential criterion for nonsingularity.
\end{proof}

\section{Formal Schemes}

\begin{definition}[Adic rings and formal spectra]
\label{def:II-9-adic}
\source{hartshorne-algebraic-geometry:II.9}
An adic ring $A$ is complete and separated for the topology defined by an ideal of definition $I$. Its
formal spectrum $\operatorname{Spf}A$ has the open prime ideals containing a power of $I$ and the completed
structure sheaf.
\end{definition}

\begin{definition}[Formal completion]
\label{def:II-9-completion}
\source{hartshorne-algebraic-geometry:II.9}
For a closed subscheme defined by $\mathcal I$, the formal completion is the ringed space represented
by the compatible thickenings $X_n=(X,\mathcal O_X/\mathcal I^{n+1})$. A coherent formal module is a
compatible inverse system of coherent modules on $X_n$.
\end{definition}

\begin{proposition}[Affine formal schemes]
\label{prop:II-9-1}
\source{hartshorne-algebraic-geometry:II.9.1}
\dcref{II.9.1}
For an adic ring $A$, sections on a basic formal open are the corresponding completed localization;
the compatible quotients recover $A/I^{n+1}$, and $\operatorname{Spf}A$ is the inverse limit of
$\operatorname{Spec}(A/I^{n+1})$.
\end{proposition}
\begin{proof}
Complete the localization along the ideal of definition. Its reductions modulo $I^{n+1}$ are the
ordinary affine structure sheaves, and completeness identifies sections with the inverse limit.
\end{proof}

\begin{proposition}[Formal functions]
\label{prop:II-9-2}
\source{hartshorne-algebraic-geometry:II.9.2}
\dcref{II.9.2}
A morphism into a formal completion is equivalent to a compatible system of morphisms into all
infinitesimal thickenings; coherent formal modules are equivalent to compatible coherent systems
under adic completeness.
\end{proposition}
\begin{proof}
Take inverse limits on affine charts. Compatibility gives a continuous map of complete rings, and
the universal property of completion gives the converse. Finite generation permits passage between
modules and their adic completions.
\end{proof}


\begin{theorem}[Completed modules]
\label{thm:II-9-3A}
\source{hartshorne-algebraic-geometry:II.9.3A}
\dcref{II.9.3A}
For a Noetherian ring $A$ and ideal $I$, completion is exact on finitely generated modules, and the
completed module is finite over the completed ring; an already complete module is recovered by completion.
\end{theorem}
\begin{proof}
Artin--Rees identifies the induced filtrations, so inverse limits preserve the exact sequence; finite
generators complete to generators, and completeness gives the canonical isomorphism.
\end{proof}

\begin{proposition}[Formal affine schemes]
\label{prop:II-9-4}
\source{hartshorne-algebraic-geometry:II.9.4}
\dcref{II.9.4}
If $X=\operatorname{Spec}A$ and $Y=V(I)$, the formal completion is
$\operatorname{Spf}\widehat A$ with reductions $\operatorname{Spec}(A/I^{n+1})$.
\end{proposition}
\begin{proof}
Basic formal opens have completed localization rings, whose reductions are the affine structure sheaves
of the thickenings; taking the inverse limit gives $\operatorname{Spf}\widehat A$.
\end{proof}

\begin{proposition}[Independence of ideal of definition]
\label{prop:II-9-5}
\source{hartshorne-algebraic-geometry:II.9.5}
\dcref{II.9.5}
Any two ideals of definition on a Noetherian formal scheme give canonically the same formal scheme and
the same category of coherent formal modules.
\end{proposition}
\begin{proof}
Powers of either ideal contain powers of the other, so the inverse systems are cofinal and have
canonically isomorphic limits; the affine identifications glue.
\end{proof}

\begin{proposition}[Coherent formal modules]
\label{prop:II-9-6}
\source{hartshorne-algebraic-geometry:II.9.6}
\dcref{II.9.6}
A coherent formal module is locally the completion of a finite module, and each reduction modulo a
power of an ideal of definition is coherent.
\end{proposition}
\begin{proof}
Use finite presentations on affine formal charts; localization and completion commute for finite modules.
\end{proof}

\begin{theorem}[Formal coherent-module equivalence]
\label{thm:II-9-7}
\source{hartshorne-algebraic-geometry:II.9.7}
\dcref{II.9.7}
Coherent formal modules are equivalent to compatible systems of coherent modules on all infinitesimal
thickenings.
\end{theorem}
\begin{proof}
Reduction gives the compatible system. Conversely, lift finite affine presentations compatibly by
completeness, glue the resulting completions, and use the affine criterion to obtain inverse functors.
\end{proof}

\begin{corollary}[Exactness of completion]
\label{cor:II-9-8}
\source{hartshorne-algebraic-geometry:II.9.8}
\dcref{II.9.8}
Completion preserves short exact sequences of coherent formal modules.
\uses{thm:II-9-7}
\end{corollary}
\begin{proof}
Apply Theorem \ref{thm:II-9-7} and exactness on each finite thickening.
\end{proof}

\begin{corollary}[Kernels, cokernels, and images]
\label{cor:II-9-9}
\source{hartshorne-algebraic-geometry:II.9.9}
\dcref{II.9.9}
Kernels, cokernels, and images of morphisms of coherent formal modules are coherent.
\uses{thm:II-9-7}
\end{corollary}
\begin{proof}
Compute these objects on every thickening and apply Theorem \ref{thm:II-9-7}.
\end{proof}
